\begin{abstract}
As memory bandwidth continues to be a major bottleneck, 
high-performance architectures such as GPUs
maximize the number of high-bandwidth memory (HBM) modules that are attached; however, as internal HBM memory bandwidth continues to increase, the GPU-to-HBM interface is becoming the bandwidth bottleneck.  
%fully utilize their die ``shoreline'' or the physical perimeter by attaching multiple high-bandwidth memory (HBM) modules. While larger reticle sizes and chiplet-based GPUs provide longer shoreline and support more HBMs, the rapid increase in internal HBM bandwidth has begun to hit the physical and packaging limits of conventional GPU-to-HBM interfaces, motivating a shift toward die-to-die (D2D) interconnects. 
In this work, to address the need for more bandwidth, we identify a new architectural opportunity from the \emph{memory} or HBM shoreline to enable additional bandwidth within a GPU system.  
%In current GPU systems, only one side of each HBM’s shoreline is used to interface with the GPU, leaving the remaining perimeter underutilized.
In particular, we propose exploiting  un-utilized memory shoreline to directly interconnect the HBMs together and create an in-package \emph{memory fabric} built on D2D interconnects. We show how the memory fabric can effectively \emph{expand} bandwidth within a GPU -- both memory bandwidth and GPU bisection bandwidth.
%by providing additional physical channels. 
We show how the memory fabric increases GPU bisection bandwidth without modifying packaging constraints and mitigates the bisection bandwidth limitations of multi-chiplet GPU architectures. In addition, since HBMs are directly interconnected, we show how the HBMs can communicate and transfer data without routing through the GPU. To support direct communication, we introduce a near-memory data-movement engine (DME) that reduces data-movement overhead by exploiting the memory fabric’s bandwidth. 
We describe the architectural changes to the HBM base die to enable the memory fabric and the software interface to enable offloading of communication to the memory fabric.  By exploiting the additional bandwidth and the HBM path diversity provided, memory-fabric based system can improve memory-intensive workloads by 42\% while accelerating MoE workloads by 36\% offloading data movement. 

%and reducing contention on GPU–to-GPU D2D links, the proposed architecture mitigates shoreline-induced congestion and substantially improves overall performance. 
\end{abstract}